\begin{center}
    {\LARGE \textbf{2017分析化学}} \\[2ex]
    {\large 时量180分钟 \quad 满分90分}
\end{center}

\vspace{0.5cm}
\noindent\textbf{注意事项:}
\begin{enumerate}[label=\arabic*., parsep=0pt, itemsep=0pt]
    \item 答题前,考生先将自己的姓名、学号、考场号、座位号填写在答题卡上,将条形码准确粘贴在考生信息条形码粘贴区。
    \item 考试结束后,将本试卷与答题纸一并收回。
    \item 回答各个问题时,将答案写在答题纸上,写在本试卷上无效。
\end{enumerate}

\vspace{0.5cm}
\noindent\textbf{一、简答题(35 pts.)}

\begin{enumerate}[label=\arabic*., itemsep=1.5ex]
    \item 目前能定量计算的影响溶解度的因素有哪些，其中哪些效应使沉淀的溶解度增大。
    \item 下列各种情况下可否获得准确分析结果？若不能，分析结果是偏高还是偏低，为什么？
    \begin{enumerate}[label=\alph*), noitemsep, topsep=0pt, partopsep=0pt]
        \item $\mathrm{pH}=4$时莫尔法滴定$\ce{Cl^-}$ (已知$\ce{H2CrO4}$的$\mathrm{p}K_{a2}=6.50$)。
        \item 佛尔哈德法滴定$\ce{Cl^-}$时，加入过量$\ce{AgNO3}$溶液后，直接用$\ce{NH4SCN}$标准溶液滴定。
        \item 法扬司法滴定$\ce{Cl^-}$时，用曙红作指示剂。
    \end{enumerate}
    \item 判断下列情况对测定结果的影响(偏高、偏低或无影响)，说明原因。
    \begin{enumerate}[label=\alph*), noitemsep, topsep=0pt, partopsep=0pt]
        \item 用$\ce{K2Cr2O7}$法测定铁矿石中铁含量时，$\ce{SnCl2}$加入过量。
        \item 用碘量法测定铜合金中的铜，在溶液温度为$60^\circ\mathrm{C}$时，加入过量的$\ce{KI}$。
    \end{enumerate}
    \item 在进行络合滴定时，为什么要加入缓冲溶液控制滴定体系保持一定的$\mathrm{pH}$值。
    \item 某弱酸$\ce{HB}$在水中的$K_a=4.2\times10^{-5}$，在水相与有机相中的分配系数$K_D=44.5$，若将$\ce{HB}$从$50.0\mathrm{mL}$水溶液中萃取到$10.0\mathrm{mL}$有机溶液中，计算$\mathrm{pH}=5.0$时的萃取率(假设$\ce{HB}$在有机相中仅以$\ce{HB}$一种型体存在)。
    \item 计算$0.10\ \mathrm{mol/L}\ \ce{H2SO4}$溶液的$\mathrm{pH}$值(已知$\ce{H2SO4}$的$K_{a2}=1.0\times10^{-2}$，一级离解完全)。
    \item 忽略离子强度的影响，计算在1，10-邻二氮菲($\ce{R}$)存在下，溶液中$\ce{H2SO4}$浓度为$1.0\mathrm{mol/L}$时，$\ce{Fe^{3+}/Fe^{2+}}$电对的条件电位。(已知在$1.0\mathrm{mol/L}\ \ce{H2SO4}$中，1，10-邻二氮菲($\ce{R}$)与$\ce{Fe^{3+}}$和$\ce{Fe^{2+}}$均能形成络合物，亚铁络合物$\ce{FeR_3^{2+}}$与高铁络合物$\ce{FeR_3^{3+}}$的稳定常数之比$K_n/K_m=2.8\times10^6$，$E^\ominus_{\ce{Fe^{3+}/Fe^{2+}}}=0.77\mathrm{V}$)
\end{enumerate}

\vspace{0.5cm}
\noindent\textbf{二、计算题(55 pts.)}

\begin{enumerate}[label=\arabic*., itemsep=1.5ex]
    \item 用$0.1000\mathrm{mol/L}\ \ce{NaOH}$滴定$0.1000\mathrm{mol/L}\ \ce{HCl}$-$0.2000\mathrm{mol/L}\ \ce{H3BO3}$的混合溶液，计算化学计量点时溶液的$\mathrm{pH}$值。若滴定终点比化学计量点高$0.50$个$\mathrm{pH}$单位，计算终点误差。(已知$\ce{H3BO3}$的$K_a=5.8\times10^{-10}$)
    \item 向含有浓度均为$0.020\mathrm{mol/L}$的$\ce{Zn^{2+}}$和$\ce{Cd^{2+}}$的待测溶液中加入$\ce{KI}$掩蔽$\ce{Cd^{2+}}$，调节溶液$\mathrm{pH}=6.0$，以$0.020\mathrm{mol/L}$的$\mathrm{EDTA}$为滴定剂进行滴定，假设终点时游离$[\ce{I^-}]=2.0\mathrm{mol/L}$，能否掩蔽$\ce{Cd^{2+}}$而滴定$\ce{Zn^{2+}}$？若能，用二甲酚橙为指示剂，终点误差为多少？（已知$\lg K_{\ce{ZnY}}=16.5$，$\lg K_{\ce{CdY}}=46.46$，$\mathrm{pH}=6.0$时$\lg\alpha_{Y(\mathrm{H})}=4.65$，$\ce{Cd^{2+}}$与$\ce{I^-}$配合物的$\lg\beta_1-\lg\beta_4$分别为$2.4$，$3.4$，$5.0$，和$6.15$，二甲酚橙指示剂变色点时$\mathrm{pZn_t}=6.5$）
    \item 计算$\ce{Cu2S}$在$\mathrm{pH}=10.0$且$\ce{NH3-NH4+}$的总浓度为$0.40\mathrm{mol/L}$溶液中的溶解度。(已知$K_{\mathrm{sp}}(\ce{Cu2S})=3.2\times10^{-49}$，$K_b(\ce{NH3})=1.8\times10^{-5}$，$\ce{H2S}$的$K_{a1}=1.3\times10^{-7}$，$K_{a2}=7.1\times10^{-15}$，$\ce{Cu^+}$与$\ce{NH3}$的配合物的$\lg\beta_1=5.93$，$\lg\beta_2=10.86$)
    \item 在$\ce{Zn^{2+} + 2Q^{2-} = ZnQ_2^{2-}}$显色反应中，当显色剂$\ce{Q^{2-}}$浓度超过$\ce{Zn^{2+}}$浓度$40$倍以上时，可认为$\ce{Zn^{2+}}$全部生成$\ce{ZnQ_2^{2-}}$，有一显色液其中$c_{\ce{Zn^{2+}}}=1.0\times10^{-3}\mathrm{mol/L}$，$c_{\ce{Q^{2-}}}=5.0\times10^{-2}\mathrm{mol/L}$，在$\lambda_{\mathrm{max}}$波长处用$1\mathrm{cm}$吸收池测得吸光度为$0.400$。在同样条件下，测得$c_{\ce{Zn^{2+}}}=1.0\times10^{-3}\mathrm{mol/L}$，$c_{\ce{Q^{2-}}}=2.5\times10^{-3}\mathrm{mol/L}$显色液时，吸光度为$0.280$，求配合物$\ce{ZnQ_2^{2-}}$的稳定常数。
    \item 为提高光度法测定微量$\ce{Pd}$的灵敏度，改用一种新的显色剂。设同一溶液，用原显色剂及新显色剂各测定$4$次，所得吸光度分别为$0.128$，$0.132$，$0.125$，$0.124$及$0.129$，$0.137$，$0.135$，$0.139$。判断新显色剂测定$\ce{Pd}$的灵敏度是否有显著提高? (置信度$95\%$)（已知$95\%$置信度时$F_{0.05,2,2}=19.00$；$F_{0.05,3,3}=9.28$；$F_{0.05,4,4}=6.39$）
    \\
    已知：
    \begin{center}
    \begin{tabular}{l|cccccc}
    $f$ & $3$ & $4$ & $5$ & $6$ & $7$ & $8$ \\
    \hline
    $t_{0.05}$ & $3.18$ & $2.78$ & $2.57$ & $2.45$ & $2.36$ & $2.31$ \\
    \hline
    $t_{0.10}$ & $2.35$ & $2.13$ & $2.02$ & $1.94$ & $1.90$ & $1.86$ \\
    \end{tabular}
    \end{center}
\end{enumerate}